\section{Related Work}
\label{sec:related-work}

Personalized agents can be organized by the capability that changes---profile,
memory, planning, or action execution---and by the mechanism that changes it
\citep{xu2026personalizedagents}.  We keep those axes separate.  A retrieved
profile can alter one response without learning a persistent policy, whereas a
trained memory writer or execution controller changes future agent behavior.
Our focus is the latter setting: offline personalized initialization followed
by auditable parameter updates from deployed evidence.

\subsection{Personalized preference and policy learning}

P-RLHF learns shared LoRA parameters together with generic and user-specific
embeddings under personalized preference objectives
\citep{li2024personalized}.  Other methods place shared structure in
low-dimensional personalized reward models \citep{liu2025sharedlora}, learned
reward bases \citep{shenfeld2025reward}, probabilistic preference bases
\citep{lee2026vrf}, or a collaborative user--response graph and mixture of
LoRA experts \citep{choi2025copl}.  C-BPO is especially close to our
cross-user reasoning: it treats other users' data as implicit negative
evidence but applies a positive--unlabeled correction to avoid penalizing
overlapping preferences \citep{ma2026cbpo}.  FedPDPO provides a federated
shared/personalized alternative \citep{zhu2026fedpdpo}.

These methods motivate the user-conditioned losses inherited by
\cref{sec:ownership-cone,sec:profile-swap}.  They do not supply
Ownership-Cone P-DPO's block-specific projection of the actual optimizer
displacement, Profile-Swap Policy Gradient's same-trajectory cross-profile
rank contrast, or the common activation-and-rollback lifecycle.  In
particular, cross-user replay in \cref{sec:ownership-cone} is a stability
constraint, not an assumption that another user's positive example is a true
negative.

\subsection{Private, modular, and continually adapted parameters}

OPPU trains one PEFT module per user \citep{tan2024oppu}; Personalized Pieces
assembles reusable fragments of shared personal adapters
\citep{tan2024pieces}; and PROPER introduces population-, group-, and
user-level stages with routed LoRA components \citep{zhang2025proper}.
Generative Adapter amortizes parameter construction in a forward pass
\citep{chen2024generativeadapter}, while Meta-LoRA learns an adaptation prior
for sparse cross-domain evidence \citep{wang2026metalora}.  These works are
natural storage, cold-start, and adaptation baselines for the shared/private
parameterization used throughout this note.

Personalized Soups trains preference-specialized policies and merges them
post hoc \citep{jang2023personalizedsoups}.  Spectral Souping similarly learns
a policy basis offline and combines outputs or parameters online without
per-user retraining \citep{chow2026spectral}.  PALM selects a compact
portfolio that approximates a family of scalarized preferences
\citep{kim2026manypreferences}.  These basis and portfolio methods are direct
alternatives to maintaining a continually updated per-user checkpoint.
SPRInG instead performs selective continual LoRA adaptation and residual
retention \citep{kim2026spring}.  Our proposals add trajectory provenance,
optimizer-state transitions, candidate loading, and rollback, and must be
compared against both the no-retraining policy-basis family and selective
continual adaptation.

\subsection{Agent decisions, memory, and active evidence}

FABLE freezes the host LLM and learns a small external per-user contextual
bandit over memory use, information acquisition, and response behavior
\citep{jin2026fable}.  It is therefore a strong execution-policy baseline:
unlike our methods, it adapts decisions without updating host-policy
parameters.  Sequential clarification work trains questions that reveal a
user's preferences \citep{montazeralghaem2025clarifying}; SELM and hybrid
preference optimization provide broader active and offline-to-online alignment
parents \citep{zhang2024selm,bose2024hybrid}.  Checkpoint-Fork Querying differs
by valuing a candidate after both exact label-conditioned optimizer branches,
rather than by label uncertainty alone.

Personal memory is also an optimization problem.  Retrieval objectives can
select or rewrite personal evidence \citep{salemi2024retrieval}, and TrustMem
learns memory-transition behavior after checking coverage, preservation, and
faithfulness \citep{yang2026trustmem}.  Muscle Memory compiles recurring user
intent into executable specialists rather than retrieving text at every
interaction \citep{omran2026muscle}.  These methods motivate the executable
state payload in Support-Handoff Replay, but none defines its irreversible
transition from stale action gradients to verified current-policy
continuations.

Finally, inference-time methods alter a completed response without necessarily
updating a persistent agent policy.  Contrasting personal-preference decoding
is one representative example \citep{bu2025cope}.  Such methods remain valid
behavioral baselines, but improvement from decoding or retrieval is not
evidence for the offline-to-online parameter-learning claims studied here.
